
\section{Localisation and spectra}
\section{The Hilbert Nullstellensatz and dimension theory}
\section{Associated primes and primary decomposition}

\begin{theorem}[Theorem 4.1]\label{mcrt-auto-theorem-4-1}\dcref{4.1}\source{matsumura-commutative-ring-theory:page-0037}\uses{}
\begin{mathematicalrecord}
(i) All the ideals of A, are of the form IA,, with I an ideal of A.
   (ii) Every prime ideal of A, is of the form pA, with p a prime ideal of
A disjoint from S, and conversely, pA, is prime in A, for every such p;
exactly the same holds for primary ideals.
\end{mathematicalrecord}
\end{theorem}
\begin{proof}
\begin{mathematicalrecord}
(i) If J is an ideal of A,, set I = .Zn A. If x = a/sEJ then
x.f(s) = f(a)~J, so that ael, and then x = (l/s)*f(a)~ZA~. The converse
inclusion IA, c J is obvious, so that J = IA,.
   (ii) If P is a prime ideal of A, and we set p = Pn A, then p is a prime
ideal of A, and from the above proof P = pA,. Moreover, since P does
not contain units of A,, we have p nS = a. Conversely, if p is a prime
ideal of A disjoint from S then
        ab
        -.-EPA,     with   s,teS=>rab~p      for some    r~s,
        s t
and since r$p we must have aEp or bEp, so that a/s or bItEpA,. One
also sees easily that l$pAs, so that pAs is a prime ideal of A,.
   For primary ideals the argument is exactly the same: if p is a primary
ideal of A disjoint from S and if rabep with reS, then since no power
of r is in p we have abEp. From this we get either u/s~pA, or (b/t)?EpAs
for some n. n
\end{mathematicalrecord}
\end{proof}


\begin{theorem}[Theorem 4.2]\label{mcrt-auto-theorem-4-2}\dcref{4.2}\source{matsumura-commutative-ring-theory:page-0038}\uses{}
\begin{mathematicalrecord}
Localisation commutes with passing to quotients by ideals.
More precisely, let A be a ring, S c A a multiplicative set, I an ideal of A
and S the image of S in A/I; then
          AdzAs = (4%
\end{mathematicalrecord}
\end{theorem}
\begin{proof}
\begin{mathematicalrecord}
Both sides have the universal property for ring homomorphisms
g:A -+ C such that
      (1) every element of S maps to a unit of C,
and (2) every element of Z maps to 0;
the isomorphism follows by the uniqueness of the solution to a universal
mapping problem. In concrete terms the isomorphism is given by
          afs mod IA,++@,        where ti=u+Z,    S=s+Z.    n

   In particular, if p is a prime ideal of A then
          A&A, 21(A/P),,.
The left-hand side is the re.sidue field of the local ring A,, whereas the
right-hand side is the field of fractions of the integral domain A/p. This
field is written K(P) and called the residue field of p.
\end{mathematicalrecord}
\end{proof}


\begin{theorem}[Theorem 4.3]\label{mcrt-auto-theorem-4-3}\dcref{4.3}\source{matsumura-commutative-ring-theory:page-0038}\uses{}
\begin{mathematicalrecord}
Let A be a ring, S c A a multiplicative             set, and
f:A -A,     the canonical map. If B is a ring, with ring homomorphisms
g:A --+B and h:B -A,      satisfying
     (1) f= 4,
and (2) for every b~Z3 there exists s~S such that g(s).kg(A),      then A,
can also be regarded as a ring of fractions of B. More precisely,
         A, = B,(,, = B,, where T= {teBlh(t)r is~- a unit of A,).
\end{mathematicalrecord}
\end{theorem}
\begin{proof}
\begin{mathematicalrecord}
We can factorise h as B -+B,        -+ A,; write CI:B=-       A, for
24       Prime ideals

the second of these maps. Now g(S) c T, so that the composite A -+ B
-+ B, factorises as A --+ A, -B,;   write fi:As -B,   for the second of
these maps. Then

so that E/J = 1, the identity map of A,. Moreover by assumption, for kB
there exist aEA and SES such that bg(s) = g(a). Hence, fl(a/s) = g(a)/g(s) =
b/l. In particular for tcT, if we take UEA, such that t/l =b(u) then
u = @p(u)= a(t/l) = h(t), so that u is a unit of A,. Hence, b/t = fl(a/s)b(u-?),
and /3 is surjective. Thus CI and fi are mutually inverse, giving an
isomorphism A, N B,. The fact that A, 2: Bqcs, can be proved similarly.
(Alternatively, this follows since T is the saturation of the multiplicative
set g(S). The reader should check this fo,r himself.) w

Corollary 1. If p is a prime ideal of A, S = A - p and B satisfies the
conditions of the theorem, then setting P = pA, n B we have A, = B,.
Proof. Under these circumstances the T in the theorem is exactly B - P.
Corollary 2. Let S c A be a multiplicative    set not containing any zero-
divisors; then A can be viewed as a subring of A,, and for any intermediate
ring A c B c A,, the ring A, is a ring of fractions of B.
Corollary 3. If S and T are two multiplicative subsets of A with S c T,
then writing T? for the image of T in A,, we have (AS)T, = A,.
Corollary 4. If S c A is a multiplicative  set and P is a prime ideal of A
disjoint from S then (AS)PA,7
                            = A,. In particular if P c Q are prime ideals of
A, then
         (A&Ap = A,.
\end{mathematicalrecord}
\end{proof}


\begin{theorem}[Theorem 4.4]\label{mcrt-auto-theorem-4-4}\dcref{4.4}\source{matsumura-commutative-ring-theory:page-0041}\uses{}
\begin{mathematicalrecord}
M, E M QA As.
\end{mathematicalrecord}
\end{theorem}
\begin{proof}
\begin{mathematicalrecord}
The map M x As -M,             defined by (m, a/s)wam/s      is A-bilinear,
so that there exists a linear map cz:M @As ----) M, such that a(m @ a/s) =
am/s. Conversely we can define 8: MS -+ M @ As by /?(m/s) = m @ (l/s);
indeed, if m/s = m?Js? then ts?m = tsm? for some t ES, and so
         m @ (l/s) = m @ (ts?ltss?) = ts?m @ (l/tss?) = tsm? @3(l/tss?)
                   = m? Q ( 1is?).
Now it is easy to check that CI and /3 are mutually inverse As-linear
maps. Hence, MS and M @AAs are isomorphic as As-modules.
\end{mathematicalrecord}
\end{proof}


\begin{theorem}[Theorem 4.5]\label{mcrt-auto-theorem-4-5}\dcref{4.5}\source{matsumura-commutative-ring-theory:page-0041}\uses{}
\begin{mathematicalrecord}
M-M,        is an exact (covariant) functor from the category
of A-modules to the category of As-modules. That is, for a morphism of
A-modules f:M - N there is a morphism fs: MS --+ N, of As-modules
such that
         (id), = id (where id is the identity map of M or Ms),
         (Sf)s = Ysfsi
 and such that an exact sequence 0 +M?-M-M?+0                   goes into an
exact sequence 0 + Mk - MS --+ Mg --t 0.
\end{mathematicalrecord}
\end{theorem}
\begin{proof}
\begin{mathematicalrecord}
To prove the exactness of 0 + M; -MS       on the last line, view M?
 as a submodule of M; then for XEM? and SES,
            x/s=0 in M,otx     =0 for some tES
                          ox/s =0 in Ms,
as required. The remaining assertions follow from the properties of the
 tensor product (see Appendix A) and from the previous theorem. (Of
course they can easily be proved directly.)    n
    It follows from this that localisation commutes with @ and with Tor,
and we will treat all this together in the section on flatness in 47.
    Let A be a ring, M an A-module and p~Spec A. There are at least
two interpretations of what it should mean that some property of A or M
holds ?locally at p?. Namely, this could mean that A, (or MJ has the
property, or it could mean that A, (or M,,) has the property for all q in
some neighbourhood U of p in Spec A. The first of these is more commonly
used, but there are cases when the two interpretations coincide. In any
case, we now prove a number of theorems which assert that a local property
implies a global one.
?4      Localisation   and Spec of a ring                                27
\end{mathematicalrecord}
\end{proof}


\begin{theorem}[Theorem 4.6]\label{mcrt-auto-theorem-4-6}\dcref{4.6}\source{matsumura-commutative-ring-theory:page-0042}\uses{mcrt-auto-theorem-1-1}
\begin{mathematicalrecord}
Let A be a ring, M an A-module and XEM. If x = 0 in M,
      for every maximal ideal p of A, then x = 0.
\end{mathematicalrecord}
\end{theorem}
\begin{proof}
\begin{mathematicalrecord}
x = 0 in M,osx = 0 for some .SEA - poann (x) + p.
    . However, if 1 $ann(x) then by Theorem 1.1, there must exist a maximal
      ideal containing ann (x). Therefore 1Eann (x), that is x = 0. H
\end{mathematicalrecord}
\end{proof}


\begin{theorem}[Theorem 4.8]\label{mcrt-auto-theorem-4-8}\dcref{4.8}\source{matsumura-commutative-ring-theory:page-0042}\uses{mcrt-auto-theorem-2-2}
\begin{mathematicalrecord}
Let A be a ring and M a finite A-module. If M OAIc(m) = 0
      for every maximal ideal m then M = 0.
\end{mathematicalrecord}
\end{theorem}
\begin{proof}
\begin{mathematicalrecord}
I      = A,/mA,,   so that M @~c(m) = M,ImM,,,,   and by NAK
      (Theorem 2.2), M 0 JC(~) = OoM,,, = 0. Thus the assertion follows from
\end{mathematicalrecord}
\end{proof}


\begin{theorem}[Theorem 4.9]\label{mcrt-auto-theorem-4-9}\dcref{4.9}\source{matsumura-commutative-ring-theory:page-0042}\uses{}
\begin{mathematicalrecord}
Let f: A --t B be a homomorphism     of rings, and M a finite
      B-modules; if M OAic(p) = 0 for every p&pec A, then M = 0.
\end{mathematicalrecord}
\end{theorem}
\begin{proof}
\begin{mathematicalrecord}
If M # 0 then by Theorem 6 there is a maximal ideal P of B such
      that M, # 0, so that by NAK, M,/PM,     # 0. If we now set p = Pn A then,
      since pM, c PM,, we have M,/pM,       # 0. Set T = B -P and S = A - p;
      then the localisation Ms = M, of M as an A-module and the localisation
      Mf,s, of M as a B-module coincide (both of them are {m/slm # M, SES} ).
      We have f(S) c T, so that
              MP = MT = O-f/&       = (Mph,
      and hence
                MP/PM,J = (M,IW,h      = Wf 0, ~P))T;
      it follows that M OA+) # 0.       w
28         Prime ideals
\end{mathematicalrecord}
\end{proof}


\begin{theorem}[Theorem 4.10]\label{mcrt-auto-theorem-4-10}\dcref{4.10}\source{matsumura-commutative-ring-theory:page-0043}\uses{mcrt-auto-theorem-2-6}
\begin{mathematicalrecord}
Let A be a ring and M a finite .4-module.
  (i) For any non-negative integer r set
           U, = {pESpec AIM, can be generated over A, by r elements};
then U, is an open subset of Spec A.
   (ii) If M is a module of finite presentation then the set
           U, = {p ESpec A IM, is a free AD-module}
is open in Spec A.
\end{mathematicalrecord}
\end{theorem}
\begin{proof}
\begin{mathematicalrecord}
(i) Suppose that A, = M,o, +. . + A,w,. Each oi is of the form
Oi = mi/si with SiEA - p and mgM, but since si is a unit of A, we can
replace Wi by mi, and so assume that oi is (the image in M, of) an element
of M. Define a linear map cp:A? --+ M from the direct sum of r copies of
A to M by (a,, . . , a,)++caiwi,  and write C for the cokernel of cp.Localising
the exact sequence A?-+ M -+ C+O at a prime ideal q, we get an
exact sequence
           A;-M,-C,+O,
and when q = p we get C, = 0. C is a quotient of M, so is finitely generated,
so that the support Supp(C) is a closed set, and hence there is an open
neighbourhood I/ of p such that C, = 0 for qE V. This means that V c U,.
(In short, if Oi,. . . , O,EM generate M, at p then they generate M, for all q
in a neighbourhood of p.)
   (ii) Suppose that M, is a free A,-module, and let wi, . . . , w, be a basis.
As in (i) there is no loss of generality in assuming that oi~M. Moreover,
if we choose a suitable D(a) as a neighbourhood of p in Spec A, ml,. . . , O,
generate M, for every qeD(a). Thus, replacing A by A, and M by M,
we can assume that the elements w 1,. . . , 0, satisfy M, = C A ,a, for every
prime ideal q of A. Then by Theorem 6,
          M/~Ao,=O,           thatisM=Aw,+...+Aw,.
(We think of replacing A by A, as shrinking Spec A down to the
neighbourhood o(a) of p.) Now, defining cp:A? -M         as above, and letting
K be its kernel, we get the exact sequence
          ,0-+K--+A?---+M+0;
moreover, K, = 0. By Theorem 2.6, K is finitely generated, so that applying
?4         Localisation    and Spec     qf a ring                                           29

(i) with   r = 0, we have that      K, = 0 for every q in a neighbourhood                V of
p; this gives (A,)? N M,,      SOthat I/ c U,.        w
\end{mathematicalrecord}
\end{proof}


\begin{theorem}[Theorem 5.1]\label{mcrt-auto-theorem-5-1}\dcref{5.1}\source{matsumura-commutative-ring-theory:page-0046}\uses{}
\begin{mathematicalrecord}
Let k be a field, L an algebraic extension of k and
al,...,  a,&      then
  (i) k[a,, . . . ,~l =kh,...,q,).
  (ii) Write q:k[X,,      . . . ,X,] --+ k(a,, . . , a,) for the homomorphism
over k which maps Xi to cci; then Ker cp is the maximal ideal generated
by n elements of the form fl(X,), f2(X1, X,),. . . ,f,,(X,, . . . ,X,), where
each .fi can be taken to be manic in Xi.
\end{mathematicalrecord}
\end{theorem}
\begin{proof}
\begin{mathematicalrecord}
Let fi(X) be the minimal polynomial of a1 over k; then (fl(X,)) is a
32         Prime ideals

maximal ideal of k[X,], so that k[aJ z k[X,]/(fl(X,)) is a field, and hence
k[~~] = @a,). Now let (p2(X) be the minimal polynomial of a2 over k(cr,);
then since k(cr,) = k[~r], the coefficients of pDz can be expressed as
polynomials in a1, and there is a polynomial f,~k[X,,  X,], manic in X,,
such that (p2(X2) = .fi(ai, X,). Thus
           kC~,>4 = 4dC~l                                                   =kt~l>Q ?v k(a,)Cxzl/(fi(crl,x,)).
Proceeding         in     the                          same way, for        l<i<n                          there     is   an
fi(X,,...,Xi)EkCXl,...,                                X,], manic in Xi, such that
           4a 1). . . ) ai] = k(a,). . . ) ai)
                            -k(cx,,...,cci-,)[xil/(.l;(crl,...,ai-,,Xi)).


Now if P(X)Ek[X,,.            . ,X,] is in the kernel of q, we have q(P) =
P(a 1,..., sr,)=O, so that P(a,,. . .,a,-r,X,,)       is divisible by f,,(c~~,. ,
CI,,-~ ,X,); dividing       P(X,, . , X,) as a polynomial           in X, by the
manic polynomial f,,(X, , . . , X,) and letting R,(X I , . . . , X,) be the remain-
der, we can write P = QJ,, + R,, with Rn(al,. . . ,cI,_ l,Xn) = 0. Similarly,
dividing R,(X, , . . . , X,) as a polynomial in X, - 1 by f,, - I(X 1, . . , X, - 1) and
letting R, - ,(X1,. . . , X,) be the remainder, we get
           R,=Q,-If,-I                            +k1,
with

           L,(a       1,...,a,-2,Xn-1,Xn)=0.
Proceeding in the same way we get P = 1 Qifi + R, with R(X,, . . . , X,) = 0;
that is R = 0 and P = 1 Qifi, so that Kercp=(f,,f,   ,..., f,,). n
   The following theorem can be regarded as a converse of Theorem 1, (i).
\end{mathematicalrecord}
\end{proof}


\begin{theorem}[Theorem 5.3]\label{mcrt-auto-theorem-5-3}\dcref{5.3}\source{matsumura-commutative-ring-theory:page-0048}\uses{}
\begin{mathematicalrecord}
Let k be a field, and let m be any maximal ideal of the
   polynomial ring k[X, , . . , X,]; then the residue class field k[X,, . . . , X,1/m
   is algebraic over k. Hence m can be generated by n elements, and in
   particular if k is algebraically closed then m is of the form m =
   (Xl-al,...,       X, - a,) for aiEk.
\end{mathematicalrecord}
\end{theorem}
\begin{proof}
\begin{mathematicalrecord}
Set k[X,, . . . , X,1/m = K, and write. cli for the image of Xi in K; then
   K = k[a, , . . . , a,]. By the previous theorem, since K is a field it is algebraic
   over k, and then by Theorem 1, (ii), m is generated by n elements. If k is
   algebraically closed then k = K, so that each Xi is congruent modulo m to
   some aiEk; then (Xi -al,...,X,-a,J               cm.     On the other hand
   (X1-%...,       X, - a,) is obviously a maximal ideal, so that equality must
   hold. a
     Let k be a field and k its algebraic closure. Suppose that CDc
   kCx i,. . . ,X,1 is a subset. An n-tuple CI= (a,, . . . , a,) of elements a,Ekis an
   algebraic zero of @ if it satisfies f(a) = 0 for every f(X)&.
\end{mathematicalrecord}
\end{proof}


\begin{theorem}[Theorem 5.4]\label{mcrt-auto-theorem-5-4}\dcref{5.4}\source{matsumura-commutative-ring-theory:page-0048}\uses{}
\begin{mathematicalrecord}
The Hilbert Nullstellensatz).
     (i) If @ is a subset of k[X l,. . . ,X,1 which does not have any algebraic
 ~ zeros then the ideal generated by @ contains 1.
         (ii) Given a subset CDof k[X,, . . . ,X,] and an element fek[X,, . . .,X,1,
      suppose that f vanishes at every algebraic zero of Q. Then some power off
  : belongs to the ideal generated by a?, that is there exist v > 0,
      BiEk[X,,     , , X,] and hi& such that f? = xgihi.
     ?Proof. (i) Let Z be the ideal generated by @; if 141 then there exists a
      maximal ideal m containing I. By the previous theorem, k[X, , . . . , X,1/m
 ?i: is algebraic over k, so that it has a k-linear isomorphic embedding 8 into
 ? k. If we set &Xi mod m) = ai then for all g(X)Em we have 0 = @g(X)) =
tj,.
34            Prime ideals

gb,     9 *. ., a,,), and
                    therefore c(= (~1~).. . , a,) is an algebraic zero of m, and
hence also of @. This contradicts the hypothesis; hence, 1~1.
   (ii) Inside k[X, , . . . , X,, Y] we consider the set 0 u (1 - Yj(X)}; then this
set has no algebraic zeros, so that by (i) it generates the ideal (1). Therefore
there exists a relation of the form
         1 = 1 Pi@, Yh(W + Q(X, VU - Yf(X)),
with h,(X)&. This is an identity in Xi,. . . , X, and Y, so that it still holds if
we substitute Y = f(X)-?. Hence we have
         l = 1 pi(x, f - ?)hi(X),
SO that multiplying    through by a suitable power of f and clearing
denominators      gives f? = c g,(X)h,(X),      with g,gk[X, , , . , X,]      and
high.        n
\end{mathematicalrecord}
\end{theorem}
\begin{proof}
\notready
\end{proof}


\begin{theorem}[Theorem 5.5]\label{mcrt-auto-theorem-5-5}\dcref{5.5}\source{matsumura-commutative-ring-theory:page-0049}\uses{}
\begin{mathematicalrecord}
Let k be a field, A a ring which is finitely generated over k, and
I a proper ideal of A; then the radical of I is the intersection of all maximal
ideals containing I, that is ,/I = nlcmm.
\end{mathematicalrecord}
\end{theorem}
\begin{proof}
\begin{mathematicalrecord}
Let A = k[a,, . . . , a,], so that A is a quotient of k[X,,.               .,X,1.
Considering the inverse image of I in k[X] reduces to the case A = k[X],
and the assertion follows from Theorem 4, (ii). n
     Compared to the result ,/I = nIcP P proved in $1, the conclusion of
Theorem 5 is much stronger. It is equivalent to the condition on a ring
that every prime ideal P should be expressible as an intersection of maximal
ideals. Rings for which this holds are called Hilbert rings or Jacobson
rings, and they have been studied independently by 0. Goldmann [l] and
W. Krull [7]. See also Kaplansky [K] and Bourbaki [B-5].
\end{mathematicalrecord}
\end{proof}


\begin{theorem}[Theorem 5.6]\label{mcrt-auto-theorem-5-6}\dcref{5.6}\source{matsumura-commutative-ring-theory:page-0049}\uses{}
\begin{mathematicalrecord}
Let k be a field and A an integral domain which is finitely
generated over k; then
             dim A = tr.deg, A.
\end{mathematicalrecord}
\end{theorem}
\begin{proof}
\begin{mathematicalrecord}
Let A = k[X 1,. . . , X,1/P, and set r = tr.deg, A. To prove that
r > dim A it is enough to show that if P and Q are prime ideals of
k[X] = k[X,, . . . , X,] with Q 3 P and Q #P then
             tr. deg, k [Xl/Q < tr. deg, k[X]/P.
The k-algebra homomorphism           k [Xl/P -+ k [Xl/Q is onto, so that tr. deg,
k(X)/Q d tr. deg,k [Xl/P is obvious. Suppose that equality holds. Let
k[X]/P = k[a,, . . . , a,] and k[X]/Q = k[pl,. . ,&,]; we can assume that
/Ii,. . . , p,. is a transcendence basis for k(j3) over k. Then a,, . . . , c(, are also
95       Hilhert   Nullstellensatz   and dimension theory                      35

    algebraically independent over k, so that they form a transcendence         basis
    for &cc) over k. Now set S = k[X, , . , X,] - { 0); S is a multiplicative   set in
    k[X]    with PnS=@        and QnS=@.          Setting R=k[X,,...,X,]          and
    K=k(X,,...,     X,), we have R,  = K[X,+   r,. . . , X,], and
             Rs/PRsNk(cr,,...,x,)Ccr,+,,...,a,l,
    ~0 that R,/PR,     is algebraic over K = k(X,,. ..,X,) 2: k(a,,.. .,r,), and
    therefore by Theorem 1, PR, is a maximal ideal of R,; but this contradicts
.   the assumptions P c Q with P # Q and Q n S = fzr.
       Now let us prove that r < dim A by induction on r. If r = 0 then, by
    Theorem 1, A is a field, so dim A = 0 and the assertion holds. Now let r > 0,
    and suppose that A = k[a,, . , r,] with g1 transcendental over k; setting
    S = k[X,] - (0) and R = k[X,, . . .,X,] we get
             R, = k(X,)[X,,    . , X,]  and R,/PR, N k(x,)[a,, . . , a,].
    Hence R,/PR,      has transcendence degree r - 1 over k(X,), so that by
    induction dim R,/PR, > r - 1. Thus there exists a strictly increasing chain
    PRs=QocQl c~~~cQ,~,ofprimeidealsofR,.IfwesetPi=QinRthen
    Pi is a prime ideal of R disjoint from S; in particular, the residue class of
    X, in R/P,-,    is not algebraic over k, and so tr.deg, R/P,-, > 0. Then
    P,-, is not a maximal ideal of R by Theorem 3, and therefore R has a
    maximal ideal P, strictly bigger than P,- r. Hence dim A = coht P 3 r. n
\end{mathematicalrecord}
\end{proof}


\begin{theorem}[Theorem 5.7]\label{mcrt-auto-theorem-5-7}\dcref{5.7}\source{matsumura-commutative-ring-theory:page-0050}\uses{}
\begin{mathematicalrecord}
Let A be a Noetherian         ring and M a finite A-module.     Set
           b(M) = sup {~(p, M) + coht pJp~Supp M};
    then M can be generated by at most b(M) elements.
       This theorem is a very important link between the number of local and
    global generators. However, there was room for improvement in the bound
    for the number of generators, and in no time R. Swan obtained a better
    bound. We will prove Swan?s bound. For this we need the concept ofj-Spec
    A introduced by Swan. This is a space having the same irreducible closed
    subsets as m-Spec A, but has the advantage of having a ?generic point?, not
    Present in m-Spec A, for every irreducible closed subset.
36        Prime ideals

    A prime ideal which can be expressed as an intersection of maximal
 ideals is called a j-prime ideal, and we write j-SpecA for the set of all
j-prime ideals. We consider j-Spec A also with its topology as a subspace
 of Spec A. Set M = m-SpecA and J = j-Spec A. If F is a closed subset of
J then there is an ideal I of A such that F = V(I)n J. One sees easily that
 a prime ideal P belongs to F if and only if P can be expressed as an
intersection of elements of F n M = V(Z)n M. Hence F is determined by
FnM, so that there is a natural one-to-one correspondence between
 closed subsets of J and of M. It follows that if M is Noetherian so is J,
and they both have the same combinatorial dimension. Now let B be an
irreducible closed subset of J, and let P be the intersection of all the
elements of B. If B = V(I)n J then I c P and we can also write
B = V(P)n J. If P is not a prime ideal then there exist f, gEA such that
f $P, gq!P and fgEP; but then
          B=(V(P+fA)nJ)u(V(P+gA)nJ),
and by definition of P there is a QEB not containing f and a Q?EB not
containing g, which implies that B is reducible, a contradiction. Therefore
P is a prime ideal. Hence PEB and B = V(P)n J. This P is called the
generic point of B. Conversely if P is any element of J then V(P)n J is
an irreducible closed subset of J, and is the closure in J of (P}. We will
write j-dim P for the combinatorial dimension of V(P)n J.
   For a finite A-module M and ~EJ we set
                          0 ifM,=O
           b(p,M) =
                         j-dim p + ~(p, M)   if M, # 0.
\end{mathematicalrecord}
\end{theorem}
\begin{proof}
\notready
\end{proof}


\begin{theorem}[Theorem 5.8]\label{mcrt-auto-theorem-5-8}\dcref{5.8}\source{matsumura-commutative-ring-theory:page-0051}\uses{}
\begin{mathematicalrecord}
Swan [l]). Let A be a ring, and suppose that m-Spec A is
a Noetherian space. Let M be a finite A-module. If
         sup {b(p, M)lp~j-Spec A} = Y < cc
then M is generated by at most Y elements.
\end{mathematicalrecord}
\end{theorem}
\begin{proof}
\begin{mathematicalrecord}
Step 1. For p&pecA     and XEM, we will say that x is basic at p if x
has non-zero image in M 0 rc(p). It is easy to see that this condition is
equivalent to ~(p, M/Ax) = ~(p, M) - 1.
\end{mathematicalrecord}
\end{proof}


\begin{theorem}[Theorem 6.1]\label{mcrt-auto-theorem-6-1}\dcref{6.1}\source{matsumura-commutative-ring-theory:page-0053}\uses{}
\begin{mathematicalrecord}
Let A be a Noetherian         ring and M a non-zero A-module.
   (i) Every maximal element of the family of ideals F = (ann (x) 10 # XE M)
is an associated prime of M, and in particular Ass(M) # a.
   (ii) The set of zero-divisors for M is the union of all the associated primes
of M.
\end{mathematicalrecord}
\end{theorem}
\begin{proof}
\begin{mathematicalrecord}
(i) We have to prove that if ann (x) is a maximal element of F then it
is prime: if a, be.4 are such that abx = 0 but bx # 0 then by maximality
ann (bx) = ann (x); hence, ax = 0.
   (ii) If ax = 0 for some x # 0 then asann (x)EF, and by (i) there is an
associated prime of M containing arm(x). n
\end{mathematicalrecord}
\end{proof}


\begin{theorem}[Theorem 6.2]\label{mcrt-auto-theorem-6-2}\dcref{6.2}\source{matsumura-commutative-ring-theory:page-0053}\uses{}
\begin{mathematicalrecord}
Let S c A be a multiplicative    set, and N an As-module.
Viewing Spec (A,) as a subset of Spec A, we have Ass,(N) = Ass,~(N). If A
is Noetherian    then for an A-module         M we have Ass(M,) ==
Ass(M)nSpec(A,).
\end{mathematicalrecord}
\end{theorem}
\begin{proof}
\begin{mathematicalrecord}
For xeN we have annA(x)=ann,,(x)nA,          so that if PeAss,dN)
then PnAeAss,(N).       Conversely if p~Ass*(N) and we choose XEN
such that p = arm,(x) then x # 0, and hence, p n S = @ and pA, is a prime
ideal of A, with pA, = ann,,(x). For the second part, if p~Ass(M)n
Spec(A,) then pnS = a, and p = arm,(x) for some XEM; if (a/s)x = 0 in
M, then there is a YES such that tax =0 in M, and t$p, taE:p gives
aep, so that ann,,(x)=pA,         and pA,gAss(M,).    Conversely, if PE
Ass(M,) then without loss of generality we have P = annA.     with XEM.
Setting p = Pn A we have P = pA,. Now p is finitely generated since A is
Noetherian, and it follows that there exists some tES such that p =
ann,(tx). Therefore pass..,.        n
\end{mathematicalrecord}
\end{proof}


\begin{theorem}[Theorem 6.3]\label{mcrt-auto-theorem-6-3}\dcref{6.3}\source{matsumura-commutative-ring-theory:page-0053}\uses{}
\begin{mathematicalrecord}
Let A be a ring and O+ M? -        M -M?   +O an exact
sequence of A-modules; then
         Ass(M) c Ass (M?) u Ass (M?).
\end{mathematicalrecord}
\end{theorem}
\begin{proof}
\begin{mathematicalrecord}
If PEASS (M) then M contains a submodule N isomorphic to A/P.
Since P is prime, for any non-zero element x of N we have arm(x) = P.
Associated primes and primary decomposition                        39
?6

Therefore if N n M? # 0 we have PEASS (M?). If N n M? = 0 then the image
of N in M? is also isomorphic to A/P, SO that PEAss(M?).      H
\end{mathematicalrecord}
\end{proof}


\begin{theorem}[Theorem 6.4]\label{mcrt-auto-theorem-6-4}\dcref{6.4}\source{matsumura-commutative-ring-theory:page-0054}\uses{}
\begin{mathematicalrecord}
Let A be a Noetherian ring and M # 0 a finite A-module.
Then there exists a chain 0 = M, c M 1 c ... c M, = M of submodules of M
such that for each i we have M,/M,- 1 N A/Pi with PiESpeC A.
\end{mathematicalrecord}
\end{theorem}
\begin{proof}
\begin{mathematicalrecord}
Choose any P,EAss(M);      then there exists a submodule M, of M
with M1 N A/P,. If M, # M and we choose any P,gAss (M/M,) then there
exists Mz c M such that M,/M, N A/P,. Continuing in the same way and
using the ascending chain condition, we eventually arrive at M, = M. n
\end{mathematicalrecord}
\end{proof}


\begin{theorem}[Theorem 6.5]\label{mcrt-auto-theorem-6-5}\dcref{6.5}\source{matsumura-commutative-ring-theory:page-0054}\uses{}
\begin{mathematicalrecord}
Let A be a Noetherian ring and M a finite A-module.
   (i) Ass(M) is a finite set.
   (ii) Ass (M) c Supp (M).
   (iii) The set of minimal elements of Ass (M) and of Supp (M) coincide.
\end{mathematicalrecord}
\end{theorem}
\begin{proof}
\begin{mathematicalrecord}
(i) follows from the previous two theorems; we need only note
that Ass (A/P) = {P}. For (ii), if 0 -+ A/P -M     is exact then SO is 0 +
A~/PA, + Mp, and therefore M, # 0. For (iii) it is enough to show that
if P is a minimal element of Supp (M) then PEASS (M). We have M, # 0
so that by Theorem 2 and (ii),
          121#Ass(M,)=Ass(M)nSpec(A,)cSupp(M)nSpec(A,)
                         = (PI.
Therefore we must have PEAss(M).          n
    Let A be a Noetherian ring and M a finite A-module. Let Pi,. . . , P, be the
minimal elements of Supp (M); then Supp (M) = V(P,) LJ?. u V(P,), and the
 v(PJ are the irreducible components of the closed set Supp(M) (see
Ex. 4.11). The prime ideals P,, . . . , P, are called the isolated associated
primes of M, and the remaining associated primes of M are called embedded
Primes. If I is an ideal of A then Supp,(A/I)       is the set of prime ideals
containing I, and the minimal prime divisors of I (that is the minimal
associated primes of the A-module A/Z) are precisely the minimal prime
ideals containing I. We have seen in Ex. 4.12 that there are only a finite
number of such primes, and Theorem 5 now gives a new proof of this.
(For examples of embedded primes see Ex. 6.6 and Ex. 8.9.)
\end{mathematicalrecord}
\end{proof}


\begin{theorem}[Theorem 6.6]\label{mcrt-auto-theorem-6-6}\dcref{6.6}\source{matsumura-commutative-ring-theory:page-0055}\uses{}
\begin{mathematicalrecord}
Let A be a Noetherian ring and M a finite A-module. Then a
submodule N c M is primary if and only if Ass (M/N) consists of one
element only. In this case, if Ass (M/N) = {P] and ann (M/N) = I then I is
primary and JI = P.
\end{mathematicalrecord}
\end{theorem}
\begin{proof}
\begin{mathematicalrecord}
If Ass(M/N) = {P> then by the previous theorem Supp(M/N) =
 V(P), so that P = ,/(ann(M/N)).       N ow if aEA is a zero-divisor for M/N it
follows from Theorem 1 that aEP, so that aEJ(ann(M/N));             hence, N is a
primary submodule of M. Conversely, if N is a primary submodule and
PgAss(M/N)        then every aeP is a zero-divisor for M/N, so that by
assumption aEJI, where I = ann (M/N). Hence P c JI, but from the
definition of associated prime we obviously have I c P, and hence JI c P,
so that P = JI. Thus Ass (M/N) has just one element JI. We prove that in
this case I is a primary ideal: let a, SEA with b$Z; if abel then ab(M/N) = 0,
but b(M/N) # 0, so that a is a zero-divisor for M/N, and therefore
aEP=JI.         w
\end{mathematicalrecord}
\end{proof}


\begin{theorem}[Theorem 6.7]\label{mcrt-auto-theorem-6-7}\dcref{6.7}\source{matsumura-commutative-ring-theory:page-0055}\uses{}
\begin{mathematicalrecord}
If N and N? are P-primary submodules of M then so is N n N?.
\end{mathematicalrecord}
\end{theorem}
\begin{proof}
\begin{mathematicalrecord}
We can embed M/(N n N?) as a submodule of (M/N) 0 (M/N?), so
that
           Ass(M/(NnN?))    cAss(M/N)uAss(M/N?)=          {P].    n

    If N c M is a submodule, we say that N is reducible if it can be written as
an intersection N = N, n N, of two submodules N,, N, with Ni # N, and
otherwise that N is irreducible; note that this has nothing to do with the
notion of irreducible modules in representation theory (= no submodules
other than 0 and M), which is a condition on M only.
    If M is a Noetherian module then any submodule N of M can be written
as a finite intersection of irreducible submodules. Proof: let 9 be the set
of submodules N c M having no such expression. If F # @ then it has a
maximal element N,. Then N, is reducible, so that N, = N, n N,, and
Ni#9.      Now each of the Ni is an intersection of a finite number of
irreducible submodules, and hence so is N,. This is a contradiction.
\end{mathematicalrecord}
\end{proof}


\begin{theorem}[Theorem 6.8]\label{mcrt-auto-theorem-6-8}\dcref{6.8}\source{matsumura-commutative-ring-theory:page-0056}\uses{}
\begin{mathematicalrecord}
Let A be a Noetherian ring and M a finite A-module.
  (i) An irreducible submodule of M is a primary submodule.
  (ii) If
             N=N,n..,nN,      with Ass (M/N,) = (Pi}
is an irredundant primary decomposition of a proper submodule N c M
then Ass(M/N) = {P,, . . . , Pl}.
   (iii) Every proper submodule N of M has a primary decomposition. If N is
a proper submodule of M and P is a minimal associated prime of M/N then
the P-primary component of N is (pp l(Np), where (pp: M -M,         is the
canonical map, and therefore it is uniquely determined by M, N and P.
\end{mathematicalrecord}
\end{theorem}
\begin{proof}
\begin{mathematicalrecord}
(i) It is enough to prove that a submodule N c M which is not
primary is reducible: replacing M by M/N we can assume that N = 0. By
Theorem 6, Ass (M) has at least two elements P, and P,. Then M contains
submodules Ki isomorphic to A/Pi for i = 1,2. Now since ann (x) = Pi for
any non-zero x~K, we must have K, n K, = 0, and hence 0 is reducible.
   (ii) We can again assume that N = 0. If 0 = N, n...n N, then M is
isomorphic to a submodule of M/N, @... 0 M/N,, so that

            ASS(M)   c ASS                  ASS (M/NJ = {PI 2.. .) Pl}?

On the other hand N,n...nN,#O,          and taking O#xEN,n...nNN,        we
have arm(x) = 0:x = N, :x. But N, : M is a primary ideal belonging to PI, SO
that Pr M c N, for some v > 0. Therefore P; x = 0; hence there exists i 3 0
such that Pix # 0 but Py l x = 0, and choosing 0 # y~P;x we have
PIY = 0. However, since YEN, n..? n N, it follows that y#N,, and by
the definition of primary submodule ann (y) c P, , so that P, = ann (y) and
42       Prime ideals

P,EAss(M).      The same works for the other Pi, and this proves that
{P1,...,P}    cAss(M).
   (iii) We have already seen thal a proper submodule has an irreducible
decomposition, so that by(i) it has a primary decomposition. Suppose that
N = N, n.-.n N, is a shortest primary decomposition, and that N, is the
P-primary component with P = P,. By Ex. 4.8 we know that N, =
V%n-        n(N,&,    and for i > 1 a power of Pi is contained in ann (M/N,);
then since Pi # P, we have (M/Ni)p = 0, and therefore (Ni), = M,. Thus
N, = (N1)P, and hence qP1(Np) = (PPI((N~)~); it is easy to check that
the right-hand side is N,. n
\end{mathematicalrecord}
\end{proof}


\begin{theorem}[Theorem 6.10]\label{mcrt-auto-theorem-6-10}\dcref{6.10}\source{matsumura-commutative-ring-theory:page-0058}\uses{}
\begin{mathematicalrecord}
If    O+M? -M         -M?    -+O is an exact sequence of
A-modules, then Att (M?) c Att (M) c Att (M?)u Att (M?).
\end{mathematicalrecord}
\end{theorem}
\begin{proof}
\begin{mathematicalrecord}
The first inclusion is trivial from the definition. For the second, let
PEAtt (M) and let N be a submodule such that M/N is P-secondary. If
M? + N = M then M/N is a non-trivial quotient of M?, hence PEAtt (M?).
If M? + N #M then M/(M? + N) is a non-trivial quotient of M? as well
as Of M/N, hence M? has a P-secondary quotient and PEAtt(M?).             H
44       Prime ideals

  An A-module M is said to be sum-irreducible           if it is neither zero nor the
sum of two proper submodules.
\end{mathematicalrecord}
\end{proof}


\begin{theorem}[Theorem 6.11]\label{mcrt-auto-theorem-6-11}\dcref{6.11}\source{matsumura-commutative-ring-theory:page-0059}\uses{mcrt-auto-theorem-6-8}
\begin{mathematicalrecord}
If M is Artinian, then it has a secondary representation.
\end{mathematicalrecord}
\end{theorem}
\begin{proof}
\begin{mathematicalrecord}
Similar to the proof of Theorem 6.8, (iii). n
   The class of modules which have secondary representations is larger
than that of Artinian modules. Sharp [8] proved that an injective module
over a Noetherian ring has a secondary representation.
\end{mathematicalrecord}
\end{proof}


\begin{theorem}[Theorem 4.7]\label{mcrt-auto-theorem-4-7}\dcref{4.7}\source{matsumura-commutative-ring-theory:page-0042}\uses{}
\begin{mathematicalrecord}
Let A be an integral domain with field of fractions K; set
      x = m-Spec A. We consider any ring of fractions of A as a subring of K.
?     Then in this sense we have

              A = n A,,,.
                   IIEX
\end{mathematicalrecord}
\end{theorem}
\begin{proof}
\begin{mathematicalrecord}
For XEK the set I= (aEAlaxEA}          is an ideal of A. Now
      XEA, is equivalent to I $ p, so that if XEA,,, for every maximal ideal
      m then 1~1, that is XEA.  w
\end{mathematicalrecord}
\end{proof}


\begin{theorem}[Theorem 5.2]\label{mcrt-auto-theorem-5-2}\dcref{5.2}\source{matsumura-commutative-ring-theory:page-0047}\uses{}
\begin{mathematicalrecord}
Let k be a field and A = k[cc 1,. . . , a,] an integral domain, and
write r = tr. deg, A for the transcendence degree of A (that is, of its field of
fractions) over k. Then if r > 0, A is not a field.
\end{mathematicalrecord}
\end{theorem}
\begin{proof}
\begin{mathematicalrecord}
Suppose that czl,, . . , q is a transcendence basis for A over k, and set
K = k(cr,, . . . , a,). Then since c(,+ 1,. . . , a, are algebraic over K, there exist
polynomials       ,fi(X,+ I ,                            , X,)EK [X,,                1,. . . , Xi], manic of degree di in Xi,
such that
            ~C~,+I~...,hJ 2.KCX,+1,...,X,ll(f,+l,...,fn)
and
            di = CKCar+ 13. . . 1ai):K(xr+    1) ?. .2 + 111.
The coefficients of fi are in K, so that for suitable 0 # gsk[a,, . . . , r,] we
have gfi~k[cr, )...) cI,][X,+1,...         ,X,1. In other words, if we set B =
4-a 1,. . . , a,, g-?1 then the fi are polynomials with coefftcients in B. We
are now going to show that A[g- ?I= B[cl,+ 1,. . , n,] is g free module over
B with basis in;=;+ 1 $?I 0 < e, < di}. Every element of B[a,+ r , . . . , a,] can
be written as P(x,.+~,. . . ,a,,) for some PcB[X,+,           , . ,X,1; dividing P as a
?5       Hilbert   Nullstellensatz   and dimension theory                       33

    polynomial in X, by f,, and replacing P by the remainder, we can assume
   that P has degree at most d, - 1 in X,; then dividing P as a polynomial
   in X,-i by f,- i, and replacing by the remainder, we can assume that P
   has degree at most d, _ 1 - 1 in X,- r. Proceeding in the same way, we
   can assume that P has degree at most di - 1 in Xi for each i; in addition,
   the elements (a110 < e < di} are linearly independent over K(cr,+ i, . . . ,
   gi- i). Hence A[g - ?1 is a free B-module. However, B is not a field; for
   Ma 1,. . . , a,] is a polynomial ring in Yvariables over k, and hence it contains
   infinitely many irreducible polynomials (the proof of this is exactly the
   same as Euclid?s proof that there exist infinitely many primes). Hence,
   there is an irreducible polynomial hek[a,, . . . , a,] which does not divide
   g, and then obviously h-?$k[a,,...         ,a,,g-I].  Therefore B contains an
   ideal Z with Z # 0, B, and since A[g-?1 is a free module over B, ZA[g-?1
   is a proper ideal of A[g-?1. Thus A[g-?1 is not a field. But if A were a
   field then we would have Al-g-?] = A, and hence A is not a field. n
\end{mathematicalrecord}
\end{proof}
